% fancy.tex
% Configuración de encabezados, pies y título de capítulo.
% Se carga una sola vez desde el preámbulo de main.tex con % fancy.tex
% Configuración de encabezados, pies y título de capítulo.
% Se carga una sola vez desde el preámbulo de main.tex con % fancy.tex
% Configuración de encabezados, pies y título de capítulo.
% Se carga una sola vez desde el preámbulo de main.tex con % fancy.tex
% Configuración de encabezados, pies y título de capítulo.
% Se carga una sola vez desde el preámbulo de main.tex con \input{fancy.tex}.
%
% Perillas rápidas:
%   - Separación de la doble línea: los dos \vskip 3pt en \headrule.
%   - Grosor de las líneas: \headrulewidth.
%   - Tamaño del encabezado: \footnotesize (cuerpo), \Large (apertura).
%   - Título del capítulo: \raggedleft (alineación), \Huge (tamaño), \vspace*/\vskip.
%   - Al cambiar un tamaño del encabezado, ajustar \headheight y \topmargin
%     según lo que pida fancyhdr en el log.

% Altura del encabezado (debe caber el \Large de la apertura de capítulo).
\setlength{\headheight}{22pt}
\addtolength{\topmargin}{-6.5pt} % compensa el \headheight para no mover el texto

% ---------------------------------------------------------------------------
% Encabezados y pies (estilo de la tesis de referencia)
% Cuerpo: sección actual en mayúsculas a la izquierda, número de página a la
% derecha y regla bajo el encabezado. Inicio de capítulo: "Capítulo N" a la
% izquierda y número de página al pie.
% ---------------------------------------------------------------------------
\pagestyle{fancy}
\fancyhf{}
\fancyhead[L]{\footnotesize\rightmark}
\fancyhead[R]{\footnotesize\thepage}
\renewcommand{\headrulewidth}{0.4pt}
\renewcommand{\footrulewidth}{0pt}

\fancypagestyle{plain}{%
  \fancyhf{}%
  \fancyhead[L]{\Large\memoriachaphead}%
  \fancyfoot[R]{\thepage}%
  \renewcommand{\headrulewidth}{0.4pt}%
  \renewcommand{\footrulewidth}{0pt}%
  % Doble línea bajo el encabezado en las páginas que abren capítulo
  % (numerado o mudo). En el resto del documento la línea es simple.
  \renewcommand{\headrule}{{%
    \hrule height \headrulewidth width \headwidth
    \vskip -\headrulewidth
    \vskip 3pt
    \hrule height \headrulewidth width \headwidth
    \vskip -\headrulewidth
    \vskip -3pt}}%
}

\makeatletter
% \memoriachaphead guarda el texto del encabezado de la página inicial de un
% capítulo numerado ("Capítulo N"); en los capítulos sin numerar queda vacío.
\newcommand{\memoriachaphead}{}
% La marca derecha guarda la sección actual (para el resto del capítulo).
\renewcommand{\chaptermark}[1]{%
  \markboth{\@chapapp\ \thechapter}{}%
  \xdef\memoriachaphead{\@chapapp\ \thechapter}}
\renewcommand{\sectionmark}[1]{%
  \markboth{\@chapapp\ \thechapter}%
          {\MakeUppercase{\ifnum\c@secnumdepth>\z@\thesection. \ \fi #1}}}
% Capítulos sin numerar (resumen, abstract, índices, bibliografía): encabezado
% vacío. La doble línea se mantiene porque igual usan el estilo `plain`.
\let\memoria@schapter\@schapter
\renewcommand{\@schapter}[1]{\gdef\memoriachaphead{}\markboth{}{}\memoria@schapter{#1}}
% El cuerpo del capítulo muestra solo el título: la etiqueta "Capítulo N"
% aparece únicamente en el encabezado de su página inicial.
\renewcommand{\@makechapterhead}[1]{%
  \vspace*{50\p@}%
  {\parindent\z@ \raggedleft \normalfont
    \interlinepenalty\@M
    \Huge\bfseries #1\par\nobreak
    \vskip 40\p@
  }}
\makeatother
.
%
% Perillas rápidas:
%   - Separación de la doble línea: los dos \vskip 3pt en \headrule.
%   - Grosor de las líneas: \headrulewidth.
%   - Tamaño del encabezado: \footnotesize (cuerpo), \Large (apertura).
%   - Título del capítulo: \raggedleft (alineación), \Huge (tamaño), \vspace*/\vskip.
%   - Al cambiar un tamaño del encabezado, ajustar \headheight y \topmargin
%     según lo que pida fancyhdr en el log.

% Altura del encabezado (debe caber el \Large de la apertura de capítulo).
\setlength{\headheight}{22pt}
\addtolength{\topmargin}{-6.5pt} % compensa el \headheight para no mover el texto

% ---------------------------------------------------------------------------
% Encabezados y pies (estilo de la tesis de referencia)
% Cuerpo: sección actual en mayúsculas a la izquierda, número de página a la
% derecha y regla bajo el encabezado. Inicio de capítulo: "Capítulo N" a la
% izquierda y número de página al pie.
% ---------------------------------------------------------------------------
\pagestyle{fancy}
\fancyhf{}
\fancyhead[L]{\footnotesize\rightmark}
\fancyhead[R]{\footnotesize\thepage}
\renewcommand{\headrulewidth}{0.4pt}
\renewcommand{\footrulewidth}{0pt}

\fancypagestyle{plain}{%
  \fancyhf{}%
  \fancyhead[L]{\Large\memoriachaphead}%
  \fancyfoot[R]{\thepage}%
  \renewcommand{\headrulewidth}{0.4pt}%
  \renewcommand{\footrulewidth}{0pt}%
  % Doble línea bajo el encabezado en las páginas que abren capítulo
  % (numerado o mudo). En el resto del documento la línea es simple.
  \renewcommand{\headrule}{{%
    \hrule height \headrulewidth width \headwidth
    \vskip -\headrulewidth
    \vskip 3pt
    \hrule height \headrulewidth width \headwidth
    \vskip -\headrulewidth
    \vskip -3pt}}%
}

\makeatletter
% \memoriachaphead guarda el texto del encabezado de la página inicial de un
% capítulo numerado ("Capítulo N"); en los capítulos sin numerar queda vacío.
\newcommand{\memoriachaphead}{}
% La marca derecha guarda la sección actual (para el resto del capítulo).
\renewcommand{\chaptermark}[1]{%
  \markboth{\@chapapp\ \thechapter}{}%
  \xdef\memoriachaphead{\@chapapp\ \thechapter}}
\renewcommand{\sectionmark}[1]{%
  \markboth{\@chapapp\ \thechapter}%
          {\MakeUppercase{\ifnum\c@secnumdepth>\z@\thesection. \ \fi #1}}}
% Capítulos sin numerar (resumen, abstract, índices, bibliografía): encabezado
% vacío. La doble línea se mantiene porque igual usan el estilo `plain`.
\let\memoria@schapter\@schapter
\renewcommand{\@schapter}[1]{\gdef\memoriachaphead{}\markboth{}{}\memoria@schapter{#1}}
% El cuerpo del capítulo muestra solo el título: la etiqueta "Capítulo N"
% aparece únicamente en el encabezado de su página inicial.
\renewcommand{\@makechapterhead}[1]{%
  \vspace*{50\p@}%
  {\parindent\z@ \raggedleft \normalfont
    \interlinepenalty\@M
    \Huge\bfseries #1\par\nobreak
    \vskip 40\p@
  }}
\makeatother
.
%
% Perillas rápidas:
%   - Separación de la doble línea: los dos \vskip 3pt en \headrule.
%   - Grosor de las líneas: \headrulewidth.
%   - Tamaño del encabezado: \footnotesize (cuerpo), \Large (apertura).
%   - Título del capítulo: \raggedleft (alineación), \Huge (tamaño), \vspace*/\vskip.
%   - Al cambiar un tamaño del encabezado, ajustar \headheight y \topmargin
%     según lo que pida fancyhdr en el log.

% Altura del encabezado (debe caber el \Large de la apertura de capítulo).
\setlength{\headheight}{22pt}
\addtolength{\topmargin}{-6.5pt} % compensa el \headheight para no mover el texto

% ---------------------------------------------------------------------------
% Encabezados y pies (estilo de la tesis de referencia)
% Cuerpo: sección actual en mayúsculas a la izquierda, número de página a la
% derecha y regla bajo el encabezado. Inicio de capítulo: "Capítulo N" a la
% izquierda y número de página al pie.
% ---------------------------------------------------------------------------
\pagestyle{fancy}
\fancyhf{}
\fancyhead[L]{\footnotesize\rightmark}
\fancyhead[R]{\footnotesize\thepage}
\renewcommand{\headrulewidth}{0.4pt}
\renewcommand{\footrulewidth}{0pt}

\fancypagestyle{plain}{%
  \fancyhf{}%
  \fancyhead[L]{\Large\memoriachaphead}%
  \fancyfoot[R]{\thepage}%
  \renewcommand{\headrulewidth}{0.4pt}%
  \renewcommand{\footrulewidth}{0pt}%
  % Doble línea bajo el encabezado en las páginas que abren capítulo
  % (numerado o mudo). En el resto del documento la línea es simple.
  \renewcommand{\headrule}{{%
    \hrule height \headrulewidth width \headwidth
    \vskip -\headrulewidth
    \vskip 3pt
    \hrule height \headrulewidth width \headwidth
    \vskip -\headrulewidth
    \vskip -3pt}}%
}

\makeatletter
% \memoriachaphead guarda el texto del encabezado de la página inicial de un
% capítulo numerado ("Capítulo N"); en los capítulos sin numerar queda vacío.
\newcommand{\memoriachaphead}{}
% La marca derecha guarda la sección actual (para el resto del capítulo).
\renewcommand{\chaptermark}[1]{%
  \markboth{\@chapapp\ \thechapter}{}%
  \xdef\memoriachaphead{\@chapapp\ \thechapter}}
\renewcommand{\sectionmark}[1]{%
  \markboth{\@chapapp\ \thechapter}%
          {\MakeUppercase{\ifnum\c@secnumdepth>\z@\thesection. \ \fi #1}}}
% Capítulos sin numerar (resumen, abstract, índices, bibliografía): encabezado
% vacío. La doble línea se mantiene porque igual usan el estilo `plain`.
\let\memoria@schapter\@schapter
\renewcommand{\@schapter}[1]{\gdef\memoriachaphead{}\markboth{}{}\memoria@schapter{#1}}
% El cuerpo del capítulo muestra solo el título: la etiqueta "Capítulo N"
% aparece únicamente en el encabezado de su página inicial.
\renewcommand{\@makechapterhead}[1]{%
  \vspace*{50\p@}%
  {\parindent\z@ \raggedleft \normalfont
    \interlinepenalty\@M
    \Huge\bfseries #1\par\nobreak
    \vskip 40\p@
  }}
\makeatother
.
%
% Perillas rápidas:
%   - Separación de la doble línea: los dos \vskip 3pt en \headrule.
%   - Grosor de las líneas: \headrulewidth.
%   - Tamaño del encabezado: \footnotesize (cuerpo), \Large (apertura).
%   - Título del capítulo: \raggedleft (alineación), \Huge (tamaño), \vspace*/\vskip.
%   - Al cambiar un tamaño del encabezado, ajustar \headheight y \topmargin
%     según lo que pida fancyhdr en el log.

% Altura del encabezado (debe caber el \Large de la apertura de capítulo).
\setlength{\headheight}{22pt}
\addtolength{\topmargin}{-6.5pt} % compensa el \headheight para no mover el texto

% ---------------------------------------------------------------------------
% Encabezados y pies (estilo de la tesis de referencia)
% Cuerpo: sección actual en mayúsculas a la izquierda, número de página a la
% derecha y regla bajo el encabezado. Inicio de capítulo: "Capítulo N" a la
% izquierda y número de página al pie.
% ---------------------------------------------------------------------------
\pagestyle{fancy}
\fancyhf{}
\fancyhead[L]{\footnotesize\rightmark}
\fancyhead[R]{\footnotesize\thepage}
\renewcommand{\headrulewidth}{0.4pt}
\renewcommand{\footrulewidth}{0pt}

\fancypagestyle{plain}{%
  \fancyhf{}%
  \fancyhead[L]{\Large\memoriachaphead}%
  \fancyfoot[R]{\thepage}%
  \renewcommand{\headrulewidth}{0.4pt}%
  \renewcommand{\footrulewidth}{0pt}%
  % Doble línea bajo el encabezado en las páginas que abren capítulo
  % (numerado o mudo). En el resto del documento la línea es simple.
  \renewcommand{\headrule}{{%
    \hrule height \headrulewidth width \headwidth
    \vskip -\headrulewidth
    \vskip 3pt
    \hrule height \headrulewidth width \headwidth
    \vskip -\headrulewidth
    \vskip -3pt}}%
}

\makeatletter
% \memoriachaphead guarda el texto del encabezado de la página inicial de un
% capítulo numerado ("Capítulo N"); en los capítulos sin numerar queda vacío.
\newcommand{\memoriachaphead}{}
% La marca derecha guarda la sección actual (para el resto del capítulo).
\renewcommand{\chaptermark}[1]{%
  \markboth{\@chapapp\ \thechapter}{}%
  \xdef\memoriachaphead{\@chapapp\ \thechapter}}
\renewcommand{\sectionmark}[1]{%
  \markboth{\@chapapp\ \thechapter}%
          {\MakeUppercase{\ifnum\c@secnumdepth>\z@\thesection. \ \fi #1}}}
% Capítulos sin numerar (resumen, abstract, índices, bibliografía): encabezado
% vacío. La doble línea se mantiene porque igual usan el estilo `plain`.
\let\memoria@schapter\@schapter
\renewcommand{\@schapter}[1]{\gdef\memoriachaphead{}\markboth{}{}\memoria@schapter{#1}}
% El cuerpo del capítulo muestra solo el título: la etiqueta "Capítulo N"
% aparece únicamente en el encabezado de su página inicial.
\renewcommand{\@makechapterhead}[1]{%
  \vspace*{50\p@}%
  {\parindent\z@ \raggedleft \normalfont
    \interlinepenalty\@M
    \Huge\bfseries #1\par\nobreak
    \vskip 40\p@
  }}
\makeatother
